%======================
% Chapter 5 : Methodology
%======================

\chapter{\scshape Methodology}

\section{System Block Diagram}

\begin{figure}[H]
	\centering
	\includegraphics[width=\textwidth]{fig_system_flow.png}
	\caption{System Block Diagram}
	\label{fig:system_flow}
\end{figure}


SPARK consists of two physical subsystems: a battery-powered wearable node and a gateway unit, connected over a local WiFi network. Figure~\ref{fig:system_flow} illustrates the top-level system architecture.

The wearable node (ESP32 DevKit V1 + MPU6050) acquires raw 6-axis IMU data at 100\,Hz, applies a complementary filter for noise rejection, and runs the two-layer fall detection pipeline entirely on-device. On a confirmed fall, a JSON event payload is published to the gateway over WiFi via MQTT. The Raspberry Pi 4B gateway ingests the event, fires a Telegram caregiver alert, updates the Streamlit dashboard, logs the event with SHAP attribution to PostgreSQL, and auto-generates a PDF clinical incident report. Figure~\ref{fig:system_flow} illustrates the end-to-end data and control flow across all pipeline stages.


Figure~\ref{fig:system_flow} is organized into four zones connected by two link types  solid arrows for sequential on-device processing, and a dashed WiFi boundary where the wearable node hands off to the gateway. Each block is described below.

\textbf{Sensor zone (wearable node, solid fill hatch)}
\begin{itemize}
	\item \textbf{MPU6050 IMU} : reads all six raw accelerometer and gyroscope axes over I$^2$C at 100\,Hz (Section~\ref{sec:imu_acquisition}).
	\item \textbf{Complementary filter} :fuses the raw axes ($\alpha = 0.98$) to reject high-frequency sensor noise before the signal reaches the detection pipeline.
	\item \textbf{Sliding window} : buffers the filtered stream in a 200-sample circular ring buffer (2 seconds at 100\,Hz), continuously overwritten until a candidate fall window is latched.
\end{itemize}

\textbf{Edge inference zone (wearable node, grid fill)}
\begin{itemize}
	\item \textbf{Two-layer detection} : Layer~1 threshold pre-filter gates Layer~2 TFLite Micro CNN inference, exactly as formalized in Algorithm 5.4.
	\item \textbf{MQTT publish} : on \texttt{CONFIRMED\_FALL}, serializes the event as JSON and publishes it to the \texttt{spark/events} topic, crossing the WiFi boundary to the gateway.
\end{itemize}

\textbf{Gateway zone (Raspberry Pi 4B, dot fill)}
\begin{itemize}
	\item \textbf{Mosquitto broker} :the local MQTT broker subscribed to \texttt{spark/events}, receiving the event published by the wearable node.
	\item \textbf{FastAPI/Uvicorn} :Incorporates the incoming event from the broker and orchestrates the remaining gateway pipeline stages.
	\item \textbf{PostgreSQL}: persists the event log, including the raw event fields and downstream SHAP output.
	\item \textbf{SHAP attribution} : computes per-feature attributions on the six peak IMU channels in under 200\,ms, identifying which channel most contributed to the fall classification.
\end{itemize}

\textbf{Output zone (parallel, dash-dot outline)}
\begin{itemize}
	\item \textbf{Telegram} : dispatches a caregiver alert within 2 seconds of event ingestion.
	\item \textbf{Streamlit Dashboard} : displays the event timestamp, severity score, and Layer~2 CNN confidence in real time.
	\item \textbf{PDF incident report} : auto-generates a one-page clinical report containing the SHAP attribution chart and a care-staff signature block.
\end{itemize}


\section{Hardware Implementation}

\subsection{Wearable Node Assembly}

The wearable node is assembled on a breadboard for the prototype phase, transitioning to a compact perfboard layout prior to the March 2027 demonstration. The assembly sequence is:

\begin{enumerate}
	\item Connect the MPU6050 to the ESP32 via I$^2$C: SDA to GPIO~21, SCL to GPIO~22, operating at 400\,kHz.
	\item Power the MPU6050 from the ESP32 3.3\,V rail.
	\item Connect the 18650 cell output (3.0--4.2\,V) to the TP4056 module (battery management), then to the AMS1117-3.3 LDO input.
	\item Connect the INA219 current sensor in series with the 3.3\,V power rail to the ESP32 VIN; connect INA219 I$^2$C lines to GPIO~21/22 (secondary address 0x41).
	\item House the assembly in a 3D-printed PLA enclosure (KEC Makerspace, Anycubic filament) with a Velcro wrist strap and flush-mounted USB-C charging port.
\end{enumerate}

\subsection{Power Budget Verification}

The INA219 provides real-time current and voltage telemetry logged to the gateway. Target metrics: active current draw $\leq$~80\,mA at 3.3\,V, providing $\geq$~25 hours of continuous monitoring from a 2000\,mAh cell. This will be measured and documented in the thesis Chapter~5 (Implementation).

\section{Firmware Implementation (ESP32 PlatformIO)}

\begin{figure}[H]
	\centering
	\includegraphics[width=0.55\textwidth]{fig_two_layer_flow.png}
	\caption{Two-layer fall detection flowchart: Layer~1 threshold pre-filter (top zone) gates Layer~2 TFLite Micro CNN inference (bottom zone). Sensitivity profile thresholds are shown in the callout.}
	\label{fig:two_layer_flow}
\end{figure}

\subsection{IMU Acquisition and Filtering}
\label{sec:imu_acquisition}

The MPU6050 is configured via I$^2$C at startup: accelerometer full-scale $\pm$8\,g, gyroscope full-scale $\pm$500$^\circ$/s, DLPF cutoff 20\,Hz. A timer interrupt service routine (ISR) fires at 100\,Hz and reads all six raw axes via direct register access (bypassing the DMP). A complementary filter ($\alpha = 0.98$) fuses accelerometer and gyroscope to produce filtered angular velocity estimates:

\begin{equation}
	\hat{\omega}_{k} = 0.98 \cdot (\hat{\omega}_{k-1} + \omega_{\text{raw},k} \cdot \Delta t) + 0.02 \cdot \theta_{\text{accel},k}
\end{equation}

Raw accelerometer values pass through without additional filtering (the DLPF suffices for noise rejection above 20\,Hz).

\subsection{Sliding Window Buffer}

A circular ring buffer stores the most recent 200 samples (2 seconds at 100\,Hz) of all six channels ($a_x, a_y, a_z, \omega_x, \omega_y, \omega_z$) in SRAM. When Layer~1 fires a \texttt{PRE\_FALL} flag, the buffer is copied to a local array, normalized using StandardScaler parameters (mean and standard deviation per channel) baked into the firmware as constant arrays at training time, and passed to the TFLite Micro interpreter.

\subsection{Layer 1 - Threshold Pre-filter}

The resultant acceleration magnitude and angular velocity magnitude are computed every ISR tick:

\begin{equation}
	|a| = \sqrt{a_x^2 + a_y^2 + a_z^2}, \qquad |\omega| = \sqrt{\omega_x^2 + \omega_y^2 + \omega_z^2}
\end{equation}

The \texttt{PRE\_FALL} flag is set when: (1) $|a|$ exceeds 2.5\,g for at least 1 sample (free-fall impact proxy), and (2) $|a|$ subsequently remains below 0.5\,g for more than 300\,ms (post-impact rest phase). The threshold values (2.5\,g, 0.5\,g, 300\,ms) will be empirically calibrated against self-collected fall simulations and are adjustable via the Strategy Pattern sensitivity profiles (Section~\ref{sec:strategy}).

\subsection{Layer 2 - TFLite Micro CNN Inference}



\section{Algorithm}

The following steps describe the workflow of the SPARK fall detection system:

\begin{enumerate}

    \item START

    \item Initialize all hardware components:
    \begin{itemize}
        \item MPU6050 IMU (6-axis accelerometer + gyroscope)
        \item ESP32 Microcontroller (520\,KB SRAM, 4\,MB Flash)
        \item WiFi module for MQTT communication
        \item Raspberry Pi 4B Gateway
        \item Mosquitto MQTT broker
    \end{itemize}

    \item Calibrate the sensor:
    \begin{itemize}
        \item Read raw accelerometer and gyroscope values at rest.
        \item Compute sensor bias/offset for calibration.
        \item Store calibration constants for runtime correction.
    \end{itemize}

    \item Continuous sensor sampling loop:
    \begin{itemize}
        \item Sample MPU6050 at 100\,Hz over I\textsuperscript{2}C.
        \item Apply complementary filter ($\alpha = 0.98$) to fuse accelerometer and gyroscope data and reject noise.
        \item Push filtered sample into a circular sliding window buffer (200 samples).
    \end{itemize}

    \item Fall detection loop:
    \begin{enumerate}
        \item Check if sliding window buffer is full (200 samples collected).
        \item Apply Layer~1 threshold check on windowed acceleration/gyroscope data.
        \item If threshold is not exceeded: discard window, return to Step~4.
        \item If threshold is exceeded: pass window to Layer~2 (1D CNN classifier).
        \item CNN outputs $P(\text{fall})$ and $P(\text{ADL})$ via softmax.
        \item If $P(\text{fall}) > 0.85$: mark event as FALL DETECTED.
        \item Else: mark as ADL (Activity of Daily Living) and return to Step~4.
        \item Log detection result with timestamp on-device.
    \end{enumerate}

    \item Event publishing and transmission:
    \begin{itemize}
        \item Package fall event as JSON payload (timestamp, confidence score, sensor snapshot).
        \item Publish payload via MQTT to topic \texttt{spark/events} over WiFi.
        \item Mosquitto broker on Raspberry Pi receives the payload ($\sim$50--100\,ms latency).
    \end{itemize}

    \item Gateway event processing:
    \begin{itemize}
        \item FastAPI/Uvicorn service ingests the incoming MQTT event.
        \item Store event record in PostgreSQL event log.
        \item Run SHAP attribution on the 6 input features ($<$200\,ms) to identify which features drove the CNN's decision.
    \end{itemize}

    \item Alert and visualization dispatch (parallel):
    \begin{itemize}
        \item Push event and SHAP explanation to Streamlit dashboard for live visualization.
        \item Send Telegram alert to caregiver/care facility ($<$2\,s).
    \end{itemize}

    \item Incident report generation:
    \begin{itemize}
        \item Compile PDF report containing: timestamp, severity, CNN confidence score, SHAP chart, and signature block.
        \item Store report and link it to the corresponding event log entry.
    \end{itemize}

    \item Return to Step~4 and continue continuous monitoring.

    \item END (on system shutdown or power-off)

\end{enumerate}



\subsection{WiFi-MQTT Event Transmission}
\label{sec:mqtt_transmission}

On \texttt{CONFIRMED\_FALL}, the ESP32 serializes the event as a JSON payload using ArduinoJSON and publishes it to the MQTT topic \texttt{spark/events} on the local Mosquitto broker. The JSON structure includes: \texttt{event\_id} (UUID), \texttt{timestamp} (ISO~8601 UTC), \texttt{layer1\_flag}, \texttt{layer2\_confidence}, \texttt{ax\_peak}, \texttt{ay\_peak}, \texttt{az\_peak}, and \texttt{severity\_score} (1--5, derived from $|a|$ magnitude). Total transmission latency from \texttt{CONFIRMED\_FALL} flag to MQTT publish is under 50\,ms over a local WiFi network.

\subsection{Sensitivity Strategy Pattern}
\label{sec:strategy}

Three sensitivity profiles are implemented as configurable parameter sets loaded from a JSON config file at gateway startup and communicated back to the ESP32 via MQTT on profile change:

\begin{itemize}
	\item \textbf{HighSensitivity} (elderly, bed-ridden): Layer~1 threshold $|a|_{\text{spike}} = 2.0$\,g, rest duration 200\,ms; Layer~2 threshold $P(\text{FALL}) > 0.75$.
	\item \textbf{Standard} (default): $|a|_{\text{spike}} = 2.5$\,g, rest duration 300\,ms; Layer~2 threshold $P(\text{FALL}) > 0.85$.
	\item \textbf{SportMode} (active ADL): $|a|_{\text{spike}} = 3.0$\,g, rest duration 400\,ms; Layer~2 threshold $P(\text{FALL}) > 0.90$.
\end{itemize}

Switching between profiles requires no firmware reflash --- the gateway updates the config file and publishes the new parameters to the ESP32 via a dedicated MQTT topic.

\section{Machine Learning Pipeline}

\subsection{Dataset Preparation}

Training data is sourced from three datasets:

\begin{itemize}
	\item \textbf{UCI MHEALTH Dataset} \cite{banos2014mhealthdroid}: 23 physical activities from 10 subjects, 3 body-worn IMUs at 50\,Hz. Resampled to 100\,Hz via linear interpolation. Fall-relevant activity windows extracted.
	\item \textbf{SisFall Dataset} \cite{sucerquia2017sisfall}: 15 fall types + 19 ADL types from 38 subjects at 200\,Hz. Downsampled to 100\,Hz. All fall types used as positive class; all ADL types used as negative class.
	\item \textbf{Self-collected KEC Dataset}: Controlled fall simulations (forward trip, lateral fall, backward fall, sitting-to-floor) and ADLs (walking, running, stair climbing, sit-to-stand, bending) recorded with the SPARK node at 100\,Hz. Target: $\geq$500 fall windows + $\geq$2000 ADL windows. This dataset will be published to GitHub as ``KEC Wearable Fall Dataset in Nepal Context.''
\end{itemize}
All datasets are windowed at 2 seconds with 50\% overlap. Per-channel StandardScaler normalization is applied; scaler parameters (mean, std per channel) are extracted and baked into the firmware as constant arrays.

\subsection{1D CNN Training}

The 1D CNN is trained in TensorFlow/Keras on the combined UCI MHEALTH + SisFall dataset, then fine-tuned on the self-collected KEC dataset using transfer learning (freeze convolutional layers, retrain Dense layers). Training protocol: Adam optimizer (lr = 1e-3), batch size 64, 50 epochs with early stopping (patience = 10 on validation loss), 80/10/10 train/validation/test split stratified by subject. Class weighting is applied to address the natural imbalance between fall events and ADL windows.

Evaluation metrics: Sensitivity (Recall for FALL class), Specificity (Recall for NON\_FALL class), F1-score, and AUC-ROC on the held-out test set. Target: Sensitivity $\geq$ 90\%, Specificity $\geq$ 90\%. Figure~\ref{fig:cnn_arch} illustrates the CNN layer sequence and the progressive compression of the time axis.

\begin{figure}[H]
	\centering
	\includegraphics[width=\textwidth]{fig_cnn_architecture.png}
	\caption{SPARK 1D CNN architecture: two Conv1D+MaxPool blocks compress the 200-sample IMU window to a 12-step feature map, which is flattened and classified by a Dense+Softmax head. INT8 quantized for TFLite Micro deployment.}
	\label{fig:cnn_arch}
\end{figure}

\subsection{INT8 Quantization and Firmware Integration}

Post-training INT8 quantization is applied via the TFLite Converter with a representative calibration dataset (200 randomly sampled training windows). The quantized TFLite model is exported as a C byte array using the \texttt{xxd -i} command and included in the PlatformIO build as a header file. Model size target: $\leq$ 120\,KB flash.

\section{Gateway Software Implementation}

\subsection{Backend and Database}

The FastAPI/Uvicorn backend exposes REST endpoints for event ingestion, profile management, and report generation. A Mosquitto MQTT consumer thread runs within the FastAPI process, subscribing to \texttt{spark/events} and persisting each \texttt{CONFIRMED\_FALL} event to the PostgreSQL \texttt{fall\_events} table (schema defined in Chapter~3, Table~3.1). The SHAP attribution pipeline (Section~\ref{sec:shap}) is invoked as a post-insert task, updating \texttt{shap\_top\_feature} and \texttt{shap\_json} columns before the Telegram alert is dispatched.

\subsection{SHAP Attribution Pipeline}
\label{sec:shap}

For each confirmed fall event, the gateway extracts the six peak IMU channel values (three accelerometer axes $a_{x}, a_{y}, a_{z}$ and three gyroscope axes $\omega_{x}, \omega_{y}, \omega_{z}$, each taken at peak magnitude) from the 2-second event window and computes SHAP attributions using the \texttt{shap} library against the trained CNN (loaded as a TFLite interpreter on the gateway). The KernelExplainer or GradientExplainer is used depending on benchmark latency on the RPi 4B. The top contributing feature (e.g., $a_{z,\text{peak}}$, contribution 42\%) is stored in \texttt{shap\_top\_feature} and the full 6-element SHAP vector is stored in \texttt{shap\_json} (JSONB). Both fields are displayed on the Streamlit dashboard per event and included in the auto-generated PDF report.

\subsection{Telegram Alert}

The python-telegram-bot library dispatches a Telegram message to the configured caregiver group within 2 seconds of \texttt{CONFIRMED\_FALL} event ingestion. The message includes: patient identifier, event timestamp (ISO~8601 local time + UTC offset), fall severity score (1--5), Layer~2 CNN confidence (as percentage), and the SHAP top-contributing IMU channel.

\subsection{Streamlit Dashboard}

The dashboard provides: (1) a live accelerometer magnitude plot (last 10 seconds, 100\,Hz WebSocket stream via FastAPI), (2) a fall event log table (most recent 50 events with timestamp, severity, confidence, SHAP top feature, and false-positive dismiss button), (3) a sensitivity profile selector (HighSensitivity / Standard / SportMode), and (4) a ``Generate Report'' button triggering PDF incident report generation for the selected event.

\subsection{PDF Incident Report}

ReportLab generates a one-page PDF clinical incident report per fall event, containing: institution header (KEC), patient identifier, event timestamp, fall severity score, Layer~2 CNN confidence score, SHAP top-contributing feature, full SHAP vector bar chart, and a signature block for care staff acknowledgement. The report format is modelled on the clinical incident reporting structure recommended by WHO \cite{who2007falls} for fall event documentation.

\section{Testing and Evaluation Plan}

\subsection{Unit Testing}

Each subsystem is validated independently before integration: (1) MPU6050 ISR correctness verified by comparing 100\,Hz output against known waveforms from a signal generator; (2) Layer~1 threshold logic verified by injecting synthetic $|a|$ sequences with known fall/non-fall labels; (3) Layer~2 TFLite Micro inference verified by comparing ESP32 INT8 outputs against Python FP32 reference on 50 held-out windows; (4) MQTT payload verified by Wireshark capture; (5) Telegram alert latency measured from MQTT publish to message receipt timestamp.

\subsection{End-to-End Latency Measurement}

A GPIO pin on the ESP32 is toggled at Layer~1 \texttt{PRE\_FALL} and at \texttt{CONFIRMED\_FALL} MQTT publish, captured on an oscilloscope to measure on-device pipeline latency. Telegram message receipt timestamp is recorded manually. Target: on-device $\leq$ 200\,ms, Telegram alert $\leq$ 2 seconds from fall event.

\subsection{ADL False-Positive Battery}

One hundred ADL sequences (walk, run, stair-climb, sit-to-stand, bend, drop object, vigorous arm swing) are performed wearing the node. The false-positive rate at Standard sensitivity profile is measured and documented. Target: false-positive rate $\leq$ 10\% on the ADL battery.

\subsection{CNN Performance Evaluation}

The held-out test set (10\% of combined dataset, stratified by subject) is used to compute Sensitivity, Specificity, F1-score, AUC-ROC, and confusion matrix. Results are reported for both the FP32 Keras model and the INT8 TFLite Micro model to quantify quantization accuracy loss.
